\chapter{Measurement details and declared approximations}
\label{app:details}
\emergencystretch=3em

This appendix holds detail the body chapters summarise: the declared approximations of the analysis,
the soundness property of the final must tier construct by construct, the check of nodes at
addresses the application does not mount, and how the figures are reproduced.

\section{Declared approximations}
\label{sec:approximations}

This section lists the constructs the analysis does not decide, so that a consumer of the graph
can tell which queries it may ask.
The list separates constructs read exactly, where a failure is a defect that Chapter~\ref{ch:results}
counts, from constructs approximated, where no completeness is claimed. Constructs of the approximated group still count in the construct denominator of Table~\ref{tab:ratios}, where they appear as dynamic, unplaced or attribution losses. A declared approximation can still cost transitions that the suites take.
The end-to-end check measures that cost on mifos (Section~\ref{sec:res-e2e}).
Two of the approximations were found by the runtime route oracle of Section~\ref{sec:res-probe} and not by
reading the analysis.
They come first because both are places where the analysis is silent rather than approximate.

A route may be declared outside the language.
alfresco's about page is registered by a JSON manifest that names a path and an identifier, which
an extension registry binds to a component class at run time.
No reader of the application's TypeScript can see such a route, and none raises a note, because
nothing in the code is unreadable.
The route is not in the code at all.
Covering it would require reading the framework's plugin format as a second source language,
which is not attempted here.

Which of several tables a prefix mounts may be decided outside the sources.
mifos holds two route tables for its users module, one per identity provider.
It selects between them with a constant assembled from \texttt{window.env}, a value injected into the
page at deployment.
Where the sources fix the condition of a conditional, the extractor reads it: a flag written in an
environment file is resolved and the arm it selects is read.
Where they do not, as here, it declares the subtree unread.
The loss is declared and is \ProbeMissedMifos{} screens wide.
Where a condition is fixed by the sources, the value read is the one the checked-out sources carry,
and a production build that substitutes another environment file may select a different arm.
The runtime route oracle reads a development build for the same reason, so the two agree about which
configuration is measured.

\textbf{Read exactly} are the following, and everything in this group is the analysis's
responsibility.
\begin{itemize}[nosep]
\item Route entries whose path is a string literal, together with the component they render.
\item \texttt{routerLink} bindings whose value is a literal or an array of literals.
\item Calls to \texttt{Router.navigate} and \texttt{Router.navigateByUrl} with a literal target.
\item The nesting relation obtained from element selectors resolved within the importing file's
scope.
\item The primary \texttt{<router-outlet>} with the shell rule that follows from it.
\item Guards named on a route, as labels on the edges entering that route's subtree.
\end{itemize}
The defects of Section~\ref{sec:res-defects} (constant route paths, lazy boundaries, unlinked nodes) lie inside this group.
They are bugs with known causes and not approximations to be declared away, which is why the cause
analysis of RQ5 can say that the final version reads them (Section~\ref{sec:res-round2}).

\textbf{Approximated} are the following constructs, each with the reason it is approximated and
the effect on the graph.

\begin{itemize}[nosep]
\item \textbf{Components created programmatically, including dialogs and overlays.} An overlay is a
user interface state that no route declares, so it has no screen node, and no edge records that a
screen opens or dismisses it. The component implementing it is reached from no template, so its sites
have the empty scope. With attribution, each such navigating component is given an unplaced origin that
carries its file, and its sites become edges from that origin, held at \emph{may}
(Section~\ref{sec:model}). A navigation in a service that only such a component injects is recorded only
as a note. The graph lacks the overlay state, the transitions into and out of it, and the screens from
which its exits fire, not the destinations of those exits.
\item \textbf{Named outlets and targets of the form \texttt{\{outlets: \ldots\}}.} An auxiliary route
adds a second active route and a second mount point, encoded in the URL as \texttt{(name:path)}
\cite{angulardocs}. The model of Chapter~\ref{ch:analysis} represents this as a second mount, and the
adapter does not implement it. A route that names a non-primary outlet, every route beneath it, and a
route whose outlet the reader cannot fix to a string mint no node.
The last kind is counted in the note that marks the route set incomplete. A commands array that carries an \texttt{\{outlets: \ldots\}}
element yields no edge and no note. In \texttt{navigateByUrl} an auxiliary group is dropped and the edge
to the primary path is held at \emph{may}, since the navigation fails when the auxiliary route does not
match and that route is not read.
A group that also carries the primary outlet yields nothing. A later
string element of a \texttt{navigate} command is one segment whatever it holds, so an auxiliary URL
written there reaches whichever route matches it as one segment. The test-source cross-check shows the
auxiliary URLs among its misses.
\item \textbf{Navigation relative to the active route (\texttt{relativeTo}).} The target is a function
of the route of the component that holds the site, and attribution is what makes it resolvable.
It is resolved once for each screen the edge is emitted from. A template \texttt{routerLink}, and a
\texttt{Router.navigate} call whose options are an object literal with a \texttt{relativeTo} key, are
resolved against that screen's path, applying a leading run of \texttt{.} and \texttt{..} segments; a
call without \texttt{relativeTo}, and \texttt{navigateByUrl}, are resolved against the root. The base is right only
when the screen is the route of the component that holds the site, and only then can the edge be must.
It is held at may where the component is rendered under a route it does not own, where \texttt{relativeTo}
is written in shorthand, passed through a variable or given another route, where a link binds its own
\texttt{[relativeTo]} input, and where the edge leaves a node without a path of its own.
Such an edge can name the wrong screen, which the may tier admits and the must tier does not.
\item \textbf{Templates passed as values (\texttt{TemplateRef}) and rendered elsewhere.} The analysis
does not follow a template value from its declaration to its instantiation. Links inside an
\texttt{<ng-template>} are read from the declaring template and attributed to the declaring component's
screens.
Every such link is held at \emph{may}, since nothing says whether or where the template is
instantiated. When the value is carried to a component rendered on other screens, the edges stay with the
declaring screens. Recovering the place of instantiation is a dataflow problem, the same obstacle as for
JavaScript XML (JSX; Section~\ref{sec:transfer}).
\item \textbf{Dispatch through NgRx effects and router event subscriptions.} The dispatch relation is
specified in Chapter~\ref{ch:analysis} and not implemented, so the screen that dispatched an action is
never recovered. At the levels that keep the call graph or its class join, \emph{full} and \emph{class},
the extractor stands in for it.
Every navigation call in a file inside the application directory that
contains a \texttt{createEffect} call is given $\top$ and offered from the shell at \emph{may}, or into
the shell's \texttt{u\_}\,sink when its target is fully dynamic. This is sound and coarse, since an
effect can fire on any screen. At the lower levels the stand-in is not applied, and an effect outside
the application directory, or declared with the older \texttt{@Effect} decorator, is treated as any
service call.
\item \textbf{Navigation inside signal effects and observable subscriptions.} The firing condition is
a data dependence in a graph the runtime builds, not a control-flow condition in the program text
(Sections~\ref{sec:signals} and~\ref{sec:reactive}). The model gives such edges the modality
\emph{unknown}, with the read set recorded as the symbolic guard.
The extractor implements neither.
A navigation inside an \texttt{effect} body, a \texttt{subscribe} callback or any other function something
else calls is attributed to the scope of its class and held at \emph{may}, because it runs only when
that function is called (Section~\ref{sec:res-round2}).
A fully dynamic target goes to the \texttt{u\_}\,sink. The graph records that the screen can be left for that destination, not under what condition.
\item \textbf{Computed navigation targets.} A target built at run time from values cannot be enumerated
statically in general, and the analysis treats it according to the static text it contains.
\begin{itemize}[nosep]
\item A target with no static part, such as a \texttt{routerLink} whose value is an interpolation among
them, is an edge of modality \emph{unknown} and confidence 0 into the \texttt{u\_}\,sink of the node it
fires from. A template literal used as an element of a commands array is read as a variable element.
\item A target with a static leading part (the head of a template literal, the left operand of a
concatenation, or the constant elements before a trailing variable element of a commands array) is
over-approximated into the \emph{may} tier.
It gets one edge to every declared route that extends that prefix by at least one segment. An empty prefix therefore connects the site to every declared route
except the root. The over-approximation is incomplete, since a prefix segment matches only an identical
route segment and a route with a parameter in that position is not a candidate.
\item A commands array whose variable element is followed by constant elements is read as a pattern and drawn as a single \emph{may} edge to the first declared route that matches it. The variable element matches only a parameter segment, so a route whose segment in that position is a constant is not reached, and no note records the omission.
\end{itemize}
\item \textbf{Injection tokens and provider bindings.} The call graph of Section~\ref{sec:implementation} resolves a call on a field through the field's declared type. Three things are not read.
\begin{itemize}[nosep]
\item A binding is taken application-wide and not from the injector that resolves the token.
So where two components provide a token differently, both implementations are joined.
This widens a scope and never narrows it.
\item A factory or value provider and a provider array a function returns from outside the application
bind nothing the pass can follow.
A call through such a token therefore contributes no edge.
It falls back to the join over every component that injects its class, or to no scope at all if none
does.
\end{itemize}
Over the corpus the pass resolved \CGDirect{} calls on application types to the class the type names,
\CGBound{} through a provider binding and \CGSubclass{} to subclasses, and left \CGUnresolved{}
unresolved, in \CGUnresolvedApps{} of the six applications.
\item \textbf{Configuration-driven route registration.} A table passed as an array to
\texttt{RouterModule.forRoot}, \texttt{forChild} or \texttt{provideRouter} is read wherever the
registration runs, including inside a provider factory the application calls.
A route
that exists only once configuration has been evaluated at startup is not recovered. Where a route object is declared and its
\texttt{path} is supplied by configuration, the path is counted as unreadable, no node is minted, and
the route set is declared a lower bound. Where the table is computed by a function and passed to a
registration call, provided through the \texttt{ROUTES} token, installed with \texttt{Router.resetConfig},
or declared in a JSON extension manifest (Section~\ref{sec:res-probe}), the routes are absent from the
graph.
No note is raised, so the route set is not declared a lower bound on their account. The
independent route parser behaves identically on each of these forms, so the two readers agree on routes
neither can see. The same approximation is why one candidate application was rejected
(Section~\ref{sec:corpus}).
\item \textbf{Guards that would let the navigation through.} A guard is a function of the application's
state, and the extractor does not evaluate it. A guard the navigation runs holds the edge at \emph{may},
even when the guard would always pass. A guard is identified by the text of its expression, so two guards that compute the same thing are two
guards. Content projected through \texttt{<ng-content>} is read as its host template places it, so a link
projected into a component that renders it only under a condition is read as unconditional. \texttt{mat-menu} and \texttt{mat-tab} are the only library containers read as rendering their content lazily. A route component behind a conditional outlet is conditional, but a component nested in it is judged by its own placement, and the application root's outlet is not read.
\item \textbf{Navigations a later one may cancel.} This item and the one on modifier keys are the
execution model of Claim~\ref{claim:must}: the claim holds only on executions they do not describe. A
navigation is demoted when code that runs after it, before the event is handled, may start another
(Section~\ref{sec:model}). Whether that code may navigate is read through the calls the reader can
resolve to a body in the application.
Such a body is a member of the class, a function of the file, or a method the checker resolves on a
typed receiver. A call into a library, through a receiver the checker cannot type,
or into code loaded at run time is taken not to navigate.
A navigation followed by such a call can therefore be published as must although the call cancels it. The reader does not follow a navigation scheduled for later, nor the other listeners one event reaches.
\item \textbf{Children of a route a function returns, mounted at the empty path.} The first empty-path
child under a screen, and the first below that one, render on the screen's own address, and their links
are drawn from it. Where the empty-path route is written inside a function that returns it, as
\texttt{withShell([\ldots])} does, the walk does not place the component the function's children name
under the screen above.
Its links stay on an unplaced origin at \emph{may}. The graph loses those
edges and asserts nothing false about them.
\item \textbf{Activation with a modifier key.} A \texttt{routerLink} on an anchor leaves a click made
with a modifier key held, or with a non-primary button, to the browser, which opens the address in a new
browsing context.
So does an anchor whose \texttt{target} names another context. The model reads every
activation of a link as a navigation of the application, so a must edge from an anchor holds for a plain
activation and not for these.
\item \textbf{The order in which Angular matches routes.} The router takes the first route, in
declaration order, that matches an address.
The matcher prefers an exact path to a parameter route and consults a redirect only when no route
matches. Wherever two declared routes, or a route and a redirect,
can answer an address, the edge is held at \emph{may} and may name the route Angular does not take. One
exception rests on an assumption. Route tables an application registers apart are joined in the order
their providers are registered, which the extractor does not read. It takes a catch-all in one such table
to be matched after every route of the others, and names each catch-all it dismissed this way in a note.
The must edges that rest on it are counted in Section~\ref{sec:res-soundness}.
\item \textbf{Ambiguous selectors in the nesting relation.} When a selector belongs to more than one
component, the candidates are narrowed to those the rendering file imports; where none is imported, the
analysis joins over all of them and records a note, which can attribute a site to screens on which its
component is never rendered.
\end{itemize}


Table~\ref{tab:approx-counts} gives the prevalence of the approximated constructs in the
three applications that were cloned for the mutation run. The in-scope kinds are not
repeated, because the construct denominator of Table~\ref{tab:conversion} already
counts them with the independent route parser. Three of the six
applications were counted, so no corpus-level prevalence is claimed.

\begin{table}[t]\small
\caption{Prevalence of the approximated constructs in three of the six applications,
produced by \texttt{approx-counts.mjs} in the harness. The figures are call-site counts
over the application sources with tests excluded, taken at the pinned corpus commits.
They are a different unit from the construct denominator of
Chapter~\ref{ch:results}, which comes from the independent route parser.}
\label{tab:approx-counts}
\TabApproxCounts
\end{table}

Two things follow from the counts.
The approximated kinds are few in realworld and bitwarden against their construct denominators
(Table~\ref{tab:conversion}: \ConstructsRealworld{} and \ConstructsBitwarden{}), so they do not explain
the recovery gap of those two.
In alfresco the approximated sites outnumber the \ConstructsAlfresco{} constructs, almost entirely
through \texttt{createEffect}, so there they cannot be set aside.
The units differ, since the denominator counts navigation constructs and the \texttt{createEffect}
figure counts effect declarations, whether or not they navigate.
The kinds are also concentrated, in different applications, so an approximation that costs nothing on
one application can be the dominant one on the next.
That is an argument for declaring the list and not reporting a single aggregate.
Call-site counts also understate the structural cost of an approximation.
A handful of \texttt{dialog.open} call sites can open many components.
The inspection of the owner components that attribution never reaches found many opened
programmatically.
The right unit is the number of components and navigation sites an approximation strands, not the number
of calls that strand them.


\section{The soundness property of the final tier}
\label{sec:res-soundness}

A gain in recovered share must be assessed together with the evidence for the emitted edges.
Table~\ref{tab:round2sound} reports the final snapshot's witness check beside RQ1's must conversion.
An edge without a detected failure is not necessarily a pass, because the checker can leave its
cited destination unread.

\begin{table}[t]\small
\caption{Final-snapshot witness-check outcomes and must conversion.
The must column counts edges whose destination was read and passed; the rest are unread
destinations, neither passes nor failures, and no must edge failed.
``Not applicable'' means the application emits no must edge. \emph{Sites cited} counts the source
locations the edges cite, \emph{Judged} counts those less the over-approximated and unread ones, and
\emph{Clean of judged} is the share of judged locations whose construct names the edge's destination.}
\label{tab:round2sound}
\TabRoundTwoSoundness
\end{table}

\subsubsection{The rules the must tier follows}

The witness check compares the destination a cited construct names with the destination its edge
names. It never asks whether the construct runs every time its code does, nor whether the router lets
the navigation through. These are the two halves of Claim~\ref{claim:must} that
Section~\ref{sec:model} assigns to the demotion rules. These rules were written during development, partly in response to
the objections of a source review and of red-team panels made with the assistant the declaration at
the front of this thesis describes. Every defect those reported was reproduced against the code
before it was repaired.

The conditions on the construct and the guards the router runs for each source and destination are those
listed with Claim~\ref{claim:must} in Section~\ref{sec:model}.

\paragraph{Same-address navigation and modality.} Angular ignores a navigation to the address its
source already displays, and where the address differs only in a parameter or the query it runs only
on some inputs. A shell or anchor edge whose destination is one of the screens it stands for is
therefore drawn per screen: must from the others, and a may self-loop at the destination. An edge
takes the modality of its most certain construct, which becomes its witness. Every other construct it
cites records the modality it earns on its own, and the witness check judges the must tier only on
the constructs that earn must. Claim~\ref{claim:must} asks that a must edge correspond to \emph{a}
navigation that occurs, not that every construct reaching the same screen does.

\subsubsection{The deterministic tier and its evidence}

Table~\ref{tab:mustledger} in Appendix~\ref{app:mustledger} records, for each construct a
development version's must tier cites, the condition the final extractor reads. It is produced with the extractor's own readers, so it records their
decisions and does not corroborate them.

With these rules the deterministic tier is \CensusMust{} edges with attribution and \CensusBaseMust{}
without it, resting on \CensusMustSites{} constructs with attribution and \CensusBaseMustSites{}
without, in \CensusMustApps{} applications. \CensusZeroMustApps{} of the \CensusApps{} have no must
edge. Most of the tier is navigation between screens that share a guarded parent, which reading guards as
the router runs them lets through. k8sdash holds \CensusMustKdash{}: the logo link and the Create button
of the toolbar its chrome component renders above every screen.

The edge count is a poor measure of how much evidence the tier holds. \CensusMustSitesKdash{}
constructs yield \CensusMustKdash{} of the \CensusMust{} edges. Without k8sdash the tier is \the\numexpr\CensusMust-\CensusMustKdash\relax{} edges on
\the\numexpr\CensusMustSites-\CensusMustSitesKdash\relax{} constructs in
\the\numexpr\CensusMustApps-1\relax{} applications, which is why RQ1 is answered in constructs. Two
assumptions weaken the evidence further.

\begin{itemize}
\item The k8sdash edges are only as right as the scope of the chrome component. An over-wide scope
there would add must edges from screens that do not render the toolbar, and over-wide scopes are not
tested outside \GTApp{} (Section~\ref{sec:res-scope}).
\item The must edges rest on an assumption about the order in which route tables registered apart
are matched (Section~\ref{sec:approximations}): \RouteOrderResting{} of them, in
\RouteOrderRestingApps{} applications (\RouteOrderRestingWhere{}). Without it the tier holds
\RouteOrderMustWithout{} edges on \RouteOrderSitesWithout{} constructs.
\end{itemize}

Each of the \CensusMustSites{} constructs was read against the source, for its condition, the address it
resolves to, and the guards its navigation runs. This reading was made once, by the AI assistant the
declaration at the front describes, and with the extractor's decisions in view. It is a check that
each decision can be defended from the source, not an independent second reading.

The \MustWitnessUnread{} unread edges are all k8s dashboard toolbar edges. These are the same edges whose source side rests on the
chrome component's scope and on the route-order assumption, so these weaknesses are not independent.
Outside k8s dashboard every must edge was read and passed. Neither the source reading nor the
mechanical check is an independent proof of Claim~\ref{claim:must}. The demotion rules are also
exercised by a mutation sweep (\SweepMutants{} mutants at \texttt{\SweepCommit}) in which every
one of the \SweepMutants{} deletions fails some test. The sweep measures whether the tests detect
deletions of the rules, not whether the rules are a complete account of runtime navigation.

\paragraph{The tier is not a property of the applications.} The extractor never evaluates a guard, so
a guard that can refuse a navigation demotes it whatever the guard would decide. Extracting the corpus
again with every guard taken to let the navigation through gives \GuardOnlyMust{} must edges in
\GuardOnlyApps{} applications instead of \CensusMust{}, resting on \GuardOnlySites{} constructs.
\GuardOnlyGuarded{} of those edges run a guard. The
small tier therefore says what the extractor can assert without evaluating guards. Evaluating them
would move at most \GuardOnlyGuarded{} edges from may to must. A reader who expected static
extraction to recover a large deterministic graph should take this as part of the answer.

\subsubsection{Over the whole graph}

Over the whole graph, the figures are counted per cited source location and not per edge, because that
is what the check opens. The \RoundTwoEdges{} edges of the final corpus graphs cite \RoundTwoSites{}
locations between them, an edge citing several wherever several constructs yield it. Three figures
answer three different questions.

\begin{itemize}
\item The must column of Table~\ref{tab:round2sound} bears on Claim~\ref{claim:must}, in one
direction: a must edge whose cited site named another destination would refute the claim. The claim's
quantification over inputs is carried by the demotion rules (Section~\ref{sec:model}).
\item The whole-graph share: only \RoundTwoGraphClean{} of the locations are clean, because
\RoundTwoOverApprox{} of them are over-approximations for which the check claims no target agreement by
design. It is reported beside the must column because \CensusMust{} edges of \CensusEdges{} are a
small slice of the graph. A precision figure for a slice is not one for the graph.
\item The judged share: of the \RoundTwoJudged{} locations the check does judge (neither
over-approximated nor among the \RoundTwoUnreadSites{} it cannot read), \RoundTwoJudgedClean{} are
clean. This is not a defect rate. The judged locations that are not clean are ones where the construct
names another destination or none, and none is cited by a must edge.
\end{itemize}

The \RoundTwoOverApprox{} figure needs a caveat, because ``by design'' reads as a property of the
method. Of the \RoundTwoEdges{} edges, \RoundTwoTopFanoutShare{} (\RoundTwoTopFanoutEdges{}) come from a
single call site in \RoundTwoTopFanoutApp{}, \RoundTwoTopFanoutSite{}, a \texttt{navigateByUrl} to a
template literal whose static prefix is empty, which over-approximates to every declared route.
The rate is therefore a property of how the method handles non-literal prefixes on this corpus, mostly
on mifos, and over-approximation remains the large majority of cited locations without that line.

\begin{table}[t]\footnotesize
\caption{The route configuration the framework holds, against the extractor's node set, as
recorded by the runtime route oracle of 29 and 30 October (Table~\ref{tab:provenance}).
\emph{Routes} counts the route entries the configuration holds once every lazy subtree is
expanded, and \emph{mounted} the addresses among them at which the router places a component on
the primary outlet. \emph{Not mounted} counts nodes the extractor declares at an address the
framework mounts nothing at. \emph{Not minted} counts addresses the framework mounts and the
extractor has no node for, and \emph{prefix only} those further addresses the framework's
configuration holds where a component is placed but no screen of the address's own is: in
\ProbeMissedPrefixUnanswered{} the route is a prefix with no empty-path child, and in
\ProbeMissedPrefixRedirect{} its empty-path child is a redirect, and in neither case is the
extractor wrong to mint no screen there.
\emph{Incomplete} is the extractor's own statement, from its soundiness notes, that part of the
route set was beyond it.}
\label{tab:probe}
\TabProbe
\end{table}


\section{Nodes at addresses the application does not mount}
\label{sec:res-unmounted}

Across the corpus the compared graph declares \ProbeInvented{} screens that the framework mounts nothing
at. Over \ProbeApps{} applications, \ProbeRoutes{} route entries and \ProbeNodes{} nodes, every
address the extractor claims is an address the framework's own configuration mounts. This is
external evidence about these six applications that no node is minted at an address the application
does not mount. It was measured on an earlier graph with the same node total (Table~\ref{tab:provenance}), which had been corrected using earlier findings of this check. It is a precondition of the must-tier property and not part of it. The
two static readers could not produce this evidence, because an address both of them place at the
wrong prefix is an address they agree on.

The evidence is bounded by what a census keyed on addresses can represent. An address is counted as
mounted when the configuration names it as a literal pattern and places a component there on the
primary outlet. \ProbeUnaddressableScreens{} route entries place a component and are counted in
neither column, because they have no such pattern: \ProbeMatcherScreens{} decides its address through
a \texttt{UrlMatcher}, and \ProbeOutletScreens{} sit on a named outlet, where the address is
\texttt{(name:path)}. Table~\ref{tab:probe-closure} shows every route entry in its bucket, and
neither table is produced unless the buckets sum to the route total and every exclusion
that drops a component-bearing entry is named. The matcher case is the more important: a navigation to
\path{/admin/users} in the running application is answered by that component, so the mounted column
is a lower bound by at least one address. How many addresses the matcher claims is not published,
because reading its body by hand is not a measurement.

The outlet entries were put to the running applications, because an exclusion that drops a component
should be shown to drop something no address reaches. Navigating the primary address each of them
sits under (\path{/personal-files} in alfresco; \path{/namespace}, \path{/event} and
\path{/clusterrole} in k8s dashboard; \path{/login}, \path{/signup} and \path{/login-with-device} in
bitwarden) is recognised as the primary screen in every case. However, the observer reads the primary
outlet's chain only, so whether an outlet component is also activated there is not measured.
Addressing the outlets in Angular's compound form, for example \path{/login(secondary:)}, produced no
recognition event from a signed-out session, which does not show that no address reaches an outlet. The exclusion therefore
rests on the ledger and not on the runtime route oracle. The alfresco viewer's entries are the ones
Section~\ref{sec:ownership} already declares outside the node set, and the rest are outside it for the
same reason, a region filled by a navigation that names the outlet. Those navigations are the construct
the ledger excludes as addressing a secondary outlet, so the two exclusions are the same fact seen from
the two ends.

\begin{table}[t]\footnotesize
\caption{Every route entry of Table~\ref{tab:probe}, in the one bucket it falls in.
\emph{One address} counts entries that a surviving entry at the same address already
answers for. \emph{Named outlet} and \emph{Matcher} are the two exclusions that drop a
component the application places (\ProbeOutletScreens{} and \ProbeMatcherScreens{}
entries, \ProbeUnaddressableScreens{} in all), and the rest place no component at any
address. This table and Table~\ref{tab:probe} are produced only if the
buckets sum to the route total and every exclusion that drops a component-bearing entry is
named here.}
\label{tab:probe-closure}
\TabProbeClosure
\end{table}

Two further limits belong with the \ProbeInvented{} figure, which is the strongest claim in this
chapter and should not be read as more than it is.

\begin{itemize}
\item ``Mounted'' is a property of the configuration and not of what a user sees. Some addresses the
extractor mints hold a component that renders only application chrome, and the address's own screen
is reached by a redirect or by an empty-path child. The runtime route oracle measures two such addresses directly:
in k8s dashboard the root redirects, and in bitwarden the root places chrome and no screen of its
own. The same shape elsewhere in the corpus is counted as a loss for one application and as
agreement for another, decided by what the extractor happened to mint there.
\item Where two routes sit at one address and a \texttt{canMatch} guard chooses between them, no
static reader can say which answers. bitwarden's \path{/vault} is such a pair: the census keeps the
first-declared component and the running router activates the other.
\end{itemize}

Neither limit moves a published number, because the census counts addresses and both members of that
pair are leaves. Both are reasons to read Table~\ref{tab:probe} as a statement about the route
configuration and not about the rendered interface.

The runtime route oracle also exposes causes that the two static readers share, which their agreement
cannot show. The first is the alfresco viewer subtree described in Section~\ref{sec:res-accounting} (history in Section~\ref{sec:named-outlet}). The second shared cause, a lazy route reached through a barrel, is described with the
other causes of the unreached route table (Section~\ref{sec:res-losses}).


\section{Reproducibility}
\label{sec:reproducibility}

The four methods that read the sources are re-run by three scripts at one commit of the extractor.
The first two refuse to start while the extractor's packages carry uncommitted changes.
\begin{itemize}
\item \texttt{regen-round.sh} runs \texttt{run-all.mjs} with the attribution pass disabled and then
enabled. It refuses an application whose checkout is not at its fixed commit. It writes a results
file with the raw output of each method. The ratio, loss, before-and-after and repair tables are
generated from those two runs.
\item \texttt{remeasure.sh} maps the corpus again and generates what is read off the graphs: the
census, the expansion check, the witness check, the guard-only run, the call-graph counts, the
ablation, the ground truth and the mutation run.
\item \texttt{mutation-sweep.mjs} records the sweep of the demotion rules apart, because it runs the
adapter's test suite once per rule.
\end{itemize}

The runtime route oracle is re-run by its own documented command, because it needs each application
built in its development configuration and served to a browser. Its result files record no extractor
commit or runtime version, so the guarantee below does not extend to them.

Each figure from the first four methods is recorded with the extractor commit, the corpus commits, the
harness version and the runtime version. Every result table is generated from those files, and prose
that states a measured figure takes it from the same file as the table. Two of the scripts that produce tables stop
rather than print a figure they cannot account for:
\begin{itemize}
\item the census stops unless every route entry falls in exactly one bucket and the buckets sum to the route total;
\item the before-and-after comparison stops unless every attribution-category construct is either recovered or classified among the remaining
losses.
\end{itemize} That ensures the accounting closes, not that each category explains the construct's cause.
Two limits remain. The evaluated commit must be one a reader can fetch, so the figures are
reproducible only once it is published. And the harness version decides several figures independently
of the extractor, so a result file is comparable only with one made by the same version.

uigraph was developed within this thesis. The version evaluated throughout is the final version:
extractor commit \texttt{c42f220}, harness 1.9.0 and mutation harness 1.9.1. Its development versions and
the corrections between them are listed in Appendix~\ref{app:repairs}.

Table~\ref{tab:provenance} gives, for each group of generated tables, the instrument and data it
comes from. The table itself, the corpus table, the counts of the rejected candidates and the
figures of superseded instruments are written by hand. Two groups are historical by design and
labelled so where printed: the rendered-interface check and the runtime route oracle. They were not
re-measured on the final version. The runtime route oracle was run on 29 and 30 October against a
graph whose extractor commit was not recorded. Its node total equals the final screen count, and that
agreement is its only link to the final graph.

\begin{table}[t]\footnotesize
\caption{Provenance of the generated tables. \emph{Final} means measured on the instrument
every other figure describes; \emph{historical} means measured on an earlier version of uigraph or its graph
and printed as a record of it.}
\label{tab:provenance}
\begin{tabularx}{\textwidth}{@{}Y>{\raggedright\arraybackslash}p{1.7cm}>{\raggedright\arraybackslash}p{3.3cm}>{\raggedright\arraybackslash}p{3.6cm}@{}}
\toprule
Tables & Status & Measured with & Regenerated from \\
\midrule
Loss ledger, before and after, repair tables, plots & final & \texttt{c42f220}, harness 1.9.0 &
\path{results-round2/final14.tgz}; the plots also \path{results-rq4/} \\
Census, must tier, witness check, guard-only, route order, call graph, expansion & final &
\texttt{c42f220} & the corpus graphs, via \texttt{remeasure.sh} \\
Ablation, structural reference & final & \texttt{c42f220}; generated at later commits with identical \texttt{packages/} &
\path{results-ablation/ablation.tgz} \\
Second structural reference (alfresco) & final & \texttt{c42f220}, one map per level &
\path{ground-truth/}\allowbreak\path{alfresco.json} \\
Mutation (RQ4) & final & \texttt{c42f220}, mutate 1.9.1 & \path{results-rq4/} \\
End-to-end navigation check & final & \texttt{c42f220} graphs (mifos \emph{rooted} mapped after scoring); suites run 3 Nov &
\path{e2e-nav/results/}, via \texttt{score.mjs} \\
Demotion-rule sweep & final & \texttt{d05227c}: \texttt{c42f220} plus the measurement switch &
\path{results-must/mutation-sweep.json} \\
Rendered-interface check (Table~\ref{tab:scope}) & historical & an earlier graph, before the call
graph; check 1.3.0 & \path{route-probe/results/} \\
Runtime route oracle (Tables~\ref{tab:probe}, \ref{tab:probe-closure}) & historical, 29--30 Oct &
oracle 1.0.0, comparison 1.1.0; extractor commit not recorded & \path{route-probe/results/} \\
Construct prevalence (Table~\ref{tab:approx-counts}) & no extractor & pinned corpus commits & \texttt{approx-counts.mjs} over the sources \\
May-reason census (Table~\ref{tab:mayreasons}) & final & \texttt{d05227c}: \texttt{c42f220} plus a measurement switch &
\path{results-may-}\allowbreak\path{reasons/} \\
Earlier work's restrictions (Table~\ref{tab:world}) & no extractor & the independent route parser's final output &
\path{results-round2/final14.tgz} \\
\bottomrule
\end{tabularx}


\end{table}

